\chapter{Выбор первой архитектуры версии V}

Задача этой главы --- выбрать одно первичное направление исследования, не
выдавая сам выбор за математическое или физическое достижение и не используя
наблюдаемые числа.

\section{Лексикографический порядок}

Вместо взвешенной оценки применяется последовательный порядок:
\begin{enumerate}
 \item не переносить действие, уже отвергнутое в версии IV;
 \item сформулировать одно совместное утверждение уровня \(N_2\);
 \item задать двоичную проверку до феноменологии;
 \item ограничить исходный класс моделей;
 \item предпочесть направление, которое устраняет общий разрыв меры и
 вакуума при наименьшем числе новых сущностей и имеет явный функционал.
\end{enumerate}

\section{Сравнение трёх классов}

Совместная вариация состояния и носителя непосредственно касается главного
разрыва меры и предлагает явный функционал, но ещё не фиксирует его смысл и
топологические веса. Граничный источник обладает сильными точными модулями,
но пока не имеет общего следа. Новая конечная геометрия допускает требуемый
знак, однако без ограничения сложности легко превращается в бесконечный
перебор.

\begin{proposition}[Выбор проверки V.0]
Первичным направлением для одной проверки выбирается совместная вариация
состояния и носителя. Граничный источник сохраняется как независимая
контрольная проверка. Новая конечная геометрия откладывается до задания
конечной области перебора.
\end{proposition}

Выбор не означает прохождения математического критерия. Число
математически и физически замкнутых архитектур остаётся равным нулю.

\section{Последствия заморозки}

Запрещено переносить коэффициенты граничного или конечно-геометрического
направления в функционал носителя; класс сравниваемых носителей фиксируется
до вычисления их порядка; функционалы Гиббса и положительного спектрального
действия нельзя складывать со свободным весом; массы, связи, \(S_{\rm vac}\)
и смешивания не могут выбирать знак или контрчлены.

\section{Результат последующих проверок}

Проверка \path{version5_carrier_measure_freeze_gate} закрыла текущую
реализацию носителя: нормировка состояния при фиксированной геометрии
проходит, но исходник не выводит меру на геометриях, единственный смысл
функционала, конечные топологические веса и полную полевую меру BV.

После этого была активирована граничная проверка
\path{version5_boundary_parent_trace_freeze_gate}; она также завершилась
отрицательно: точная пара вильсоновских коэффициентов существует только на
операторном уровне, а единый след, выбор заряда или источника и совместный
вывод конденсата, оси и моды Майораны отсутствуют.

Оставшийся класс новой конечной геометрии допускается только после проверки
\path{version5_finite_geometry_complexity_bound_gate}.

\section{Вывод}

\begin{protocol}
\begin{enumerate}
 \item первичная проверка: \textbf{состояние и носитель};
 \item независимая контрольная проверка: \textbf{граничный источник};
 \item новая конечная геометрия: \textbf{отложена до ограничения сложности};
 \item математическое замыкание: \textbf{не достигнуто};
 \item физическое замыкание: \textbf{не достигнуто}.
\end{enumerate}
\end{protocol}

Машинный реестр сохранён в
\path{s2t_v5_architecture_selection_gate_results.json}.